\section{Introduction}
This year marks the formation of a new team from the University of Augsburg participating in the RoboCup@Home competition. We work in collaboration with Aracom, whose expertise supported the team throughout the development process. Our team is comprised of Robert Jeutter from Aracom, Nils Mandischer from UniA, and the eight students in the Autonomous Robotics course. While the team is newly founded, it builds upon the technical expertise and academic background of our institution, aiming to contribute innovative approaches to robot perception, manipulation, and human–robot interaction.

For our first participation, we aim to establish a stable foundation for autonomous service robotics research. This means that we are not aiming for immediate top-level performance, but rather prioritize software architecture, maintainability, and usability. This foundation is intended to be further expanded by future students to steadily increase the robot's capabilities in the coming semesters.

We intend to lay the groundwork for long-term success in RoboCup@Home, both in terms of competition performance and in providing a valuable research and teaching environment at the University of Augsburg.
